\section{对称操作引出群}
\label{sec:04-Discrete-Mathematics/06-Groups-and-Symmetry/01}

你面前有一个正方形，卡片大小，四个角没有做任何标记——它就是一张完全对称的正方形。

把它绕中心旋转 \(90^\circ\)，放回桌面。方框的轮廓和之前完全重叠，看不出任何区别。沿竖直中线翻折，放回去，还是一模一样。再试试：先翻折再旋转 \(90^\circ\)——嗯，轮廓仍然重叠，但这个组合结果能不能用单独一次翻折或旋转来等效？

这就引出了两个方向的问题。列举方向：到底有多少种不同的操作能让正方形``看起来不变''？结构方向：这些操作之间有什么关系？它们能复合吗？复合的次序重要吗？能不能选出少数几个基本操作，让其余的操作都由它们组合产生？

动手做比盯着看更有用。拿一张正方形纸片，在背面标上四个角 A、B、C、D，然后开始操作它。做完每一个动作后，记下各个角的新位置。很快你会发现做了某些组合之后，角的位置分布和你之前记录的某个单独操作完全相同——翻折一次再旋转 \(90^\circ\)，跟你直接做另一种翻折，效果是一样的。

群不研究正方形好不好看，群研究的是可逆操作为什么能复合，以及复合的次序本身携带了哪些信息。

\subsection{从正方形的操作开始}

记逆时针旋转 \(90^\circ\) 为 \(r\)，沿过中心的一条竖直线翻折为 \(s\)。什么也不做记为 \(e\)（单位元）。

连续旋转自然产生：
\[
r^2,\quad r^3,\quad r^4=e.
\]
四次 \(90^\circ\) 旋转等于回到原位——旋转的阶是 4。

翻折两次也回到原位：
\[
s^2=e.
\]
翻折的阶是 2。

现在把旋转和翻折混在一起。先做一次翻折 \(s\)，再做一次旋转 \(r\)，记为 \(rs\)。这里约定乘积 \(ab\) 表示先做 \(b\)，再做 \(a\)，与函数复合 \((a\circ b)(x)=a(b(x))\) 的方向一致。不同教材可能从左向右读操作；只要一方约定清楚，所有结论等价，但在同一推导中不能中途换方向。

正方形的全部对称操作可以写成这八个：
\[
e,\; r,\; r^2,\; r^3,\; s,\; sr,\; sr^2,\; sr^3.
\]
注意 \(rs\) 不在这个列表里——它等于 \(sr^3\)。为什么？动手验证：先旋转 \(90^\circ\) 再翻折，和先翻折再反向旋转 \(90^\circ\)，结果相同。写成等式：
\[
rs = sr^{-1} = sr^3.
\]

这八个操作，配上一个二元运算（先做一个、再做另一个），表现出四项共同性质：
\begin{itemize}
\tightlist
\item
  任意两个对称操作的连续执行，结果仍在这八个操作之中——不会跑出去；
\item
  三个操作无论先把前两个还是后两个看成一组，最终结果相同；
\item
  存在``什么也不做''的操作 \(e\)，加上它不改变任何操作；
\item
  每个操作都有一个撤销动作：旋转的撤销是反向旋转，翻折的撤销是再翻折一次。
\end{itemize}

这四项概括了``操作可以稳定复合且每步可逆''所需的最低条件，每一项都有不可替代的具体作用。

\subsection{群的定义：四条公理}

一个群是一个集合 \(G\) 连同二元运算 \((a,b)\mapsto ab\)，满足：

\begin{enumerate}
\tightlist
\item
  封闭性：若 \(a,b\in G\)，则 \(ab\in G\)；
\item
  结合律：\((ab)c=a(bc)\)；
\item
  单位元：存在 \(e\in G\)，使 \(ea=ae=a\) 对所有 \(a\in G\) 成立；
\item
  逆元：每个 \(a\in G\) 都有 \(a^{-1}\in G\)，使 \(aa^{-1}=a^{-1}a=e\)。
\end{enumerate}

若还额外满足 \(ab=ba\) 对所有元素成立，称为交换群或 Abel 群。正方形的对称群不满足交换律——\(rs\neq sr\)，所以群公理不要求交换律。

结合律与交换律容易搞混。结合律只动括号，不动元素顺序：
\[
(ab)c = a(bc).
\]
它保证一长串操作 \(a_1a_2\cdots a_n\) 不必反复声明括号——任何加括号方式的结果都相同。交换律则是真的交换顺序 \(ab=ba\)，强得多，也经常失败。

\subsection{为什么四条都需要：少一条会怎样}

如果去掉封闭性，两步合法操作的组合可能跑出研究对象。正整数在减法下不封闭：\(2-3=-1\) 不是正整数。你不能在正整数里自由地做减法然后指望结果还在。

如果去掉结合律，``连续做三个操作''会依赖你计算时的括号位置。减法满足正整数不满足的性质，但在整数上它满足封闭性，结合律却失败了：
\[
(5-3)-1=1,\qquad 5-(3-1)=3.
\]
同一串符号，不同的括号给出不同的结果。没有结合律，你甚至不能把三个操作写成 \(a_1a_2a_3\)——必须声明是先合并前两个还是后两个。

如果去掉单位元，就没有``什么都不做''的参照。如果去掉逆元，每个过程都不可撤销——往前走可以，但永远回不到起点。

整数在乘法下：有单位元 1，有封闭性（整数乘整数仍是整数），乘法也结合。但 2 的倒数 \(1/2\) 不是整数——缺少逆元。所以整数在乘法下不是群。非零实数在乘法下才是群：每个非零实数 \(x\) 的倒数 \(1/x\) 仍是非零实数。

可逆性也划出了群的边界。把文件永久删除、把数据做有损压缩、把向量投影到低维子空间——这些操作丢掉了信息，不能撤销。它们的集合在合适的运算下通常构成幺半群（有单位元的半群），而不是群。取矩阵的秩、求两个数的最大值——这些运算甚至不满足消去律。看见一个二元运算不能直接叫群——四条公理，逐条核查，少一条都不行。

这个核查过程本身也在训练一种思维：当我们说一个操作集合``具有良好的结构''时，具体是在说哪几条性质？每一条性质能推导出什么事实？去掉某一条，哪些推理链会断掉？群论的力量不只在于定义了群，而在于它让你能精确回答``这个系统缺少什么所以不能做某件事''。

\subsection{由四条公理推出的第一批事实}

四条公理看似简单，但它们是种子——下面这些事实不是外加的，是从公理推导出来的。

\textbf{单位元唯一。}若 \(e\) 与 \(e'\) 都是单位元，则
\[
e = ee' = e'.
\]
第一个等号用了``\(e'\) 是右单位元''，第二个等号用了``\(e\) 是左单位元''。左右夹击，证明唯一。

\textbf{逆元唯一。}若 \(b,c\) 都是 \(a\) 的逆元：
\[
b = be = b(ac) = (ba)c = ec = c.
\]
这里每一步都在用公理：单位元的定义、已知 \(ac=e\)、结合律、已知 \(ba=e\)、单位元的定义。没有哪一步是凭空跳出来的。

\textbf{群中可消去。}若 \(ab=ac\)，左乘 \(a^{-1}\)：
\[
a^{-1}(ab) = a^{-1}(ac) \;\Longrightarrow\; (a^{-1}a)b = (a^{-1}a)c \;\Longrightarrow\; eb = ec \;\Longrightarrow\; b = c.
\]
但消去方向必须对。在非交换群中，从 \(ba=ca\) 应该右乘 \(a^{-1}\)。逆元不能随意搬到另一侧——左右交换可能改变结果。

\textbf{乘积的逆元会反转顺序（socks-shoes 法则）。}
\[
(ab)^{-1}=b^{-1}a^{-1},
\]
因为
\[
(ab)(b^{-1}a^{-1}) = a(bb^{-1})a^{-1} = aea^{-1} = aa^{-1} = e.
\]
生活的直觉完全一致：早上先穿袜子再穿鞋，晚上必须先脱鞋再脱袜子。顺序不一样，只会把袜子扯破。

\textbf{线性方程在群中的解唯一。}若 \(ax=b\)，则 \(x=a^{-1}b\)，且这个解是唯一的。因为如果 \(ax_1=ax_2\)，左乘 \(a^{-1}\) 即得 \(x_1=x_2\)。同样的逻辑给出 \(xa=b\) 的唯一解为 \(x=ba^{-1}\)。在非交换群中，这两个方程的解一般不同，因为 \(a^{-1}b\neq ba^{-1}\) 除非 \(a\) 与 \(b\) 交换。群中的方程求解比实数方程多一层警觉：搞清楚未知数在已知元素的左边还是右边。

这些从公理推出来的小事实，每一个都对应着你在操作正方形时的那份直觉——``撤销操作是唯一的''、``方程的解是唯一的''、``先穿袜后穿鞋，脱的时候顺序反着来''。抽象公理和手上操作之间不需要什么桥梁，因为它们说的是同一件事。

\subsection{循环群：一个操作反复生成全部元素}

若某个元素 \(g\) 的整数次幂能走遍群中全部元素，
\[
G=\set{g^k:k\in\Z},
\]
称 \(G\) 为循环群，\(g\) 为生成元。

最熟悉的例子：模 \(n\) 加法群
\[
\Z_n=\set{0,1,\ldots,n-1}
\]
由 \(1\) 生成。群运算是``加法后取模''：
\[
0,\;1,\;2,\;\ldots,\;n-1,\;0,\;\ldots
\]
每次加 1，循环出现。单位元是 0，\(a\) 的逆元是 \(-a\bmod n\)（即 \(n-a\) 当 \(a\neq0\)，0 的逆元是 0）。

元素 \(g\) 的\textbf{阶}是最小正整数 \(m\) 使 \(g^m=e\)；若不存在，阶为无穷。正方形旋转 \(r\) 的阶为 4（四次 \(90^\circ\) 回到原位），翻折 \(s\) 的阶为 2（翻两次回到原位）。阶描述的是``重复多少次回到起点''——它是元素的循环节奏。

所有循环群都是交换群：
\[
g^a g^b = g^{a+b} = g^{b+a} = g^b g^a.
\]
反过来不成立：有些交换群不能由单个元素生成。例如 Klein 四元群 \(\Z_2\times\Z_2\)（四个元素，每个非单位元素重复两次就回到单位元）。它的元素用数对表示：\((0,0),(1,0),(0,1),(1,1)\)，运算为逐分量模 2 加法。每个非零元素的阶都是 2：\((1,0)+(1,0)=(0,0)\)，\((0,1)+(0,1)=(0,0)\)，\((1,1)+(1,1)=(0,0)\)。它们的生成能力最多走遍自己和单位元，不可能独自走遍全部四个元素。所以这个群是交换群但不是循环群——交换性和单生成元性是两个独立的性质。

一个有限循环群中，\(g^k\) 的阶等于 \(m/\gcd(m,k)\)，其中 \(m\) 是 \(g\) 的阶。这意味着模 \(n\) 加法群 \(\Z_n\) 中，元素 \(k\) 生成整个群当且仅当 \(\gcd(k,n)=1\)。例如 \(\Z_{12}\) 中，生成元是 1、5、7、11——恰好是 12 的互素数。这个事实把群论又重新拉回到数论：生成元的数量是 Euler totient 函数 \(\varphi(n)\)。看似不相干的领域在循环群的阶结构这里交汇了一下。

\subsection{置换群：重排本身就是操作}

集合 \(\{1,\ldots,n\}\) 上的所有双射（一一对应）称为置换。置换在函数复合下构成\textbf{对称群} \(S_n\)：
\begin{itemize}
\tightlist
\item
  双射的复合仍是双射——封闭；
\item
  函数复合天然满足结合律——结合；
\item
  恒等映射 \(x\mapsto x\) 是单位元——什么都不动；
\item
  每个双射都有逆映射——可撤销。
\end{itemize}

\(S_n\) 有 \(n!\) 个元素。\(S_3\) 有 6 个置换，用轮换记号写：
\[
e,\;(12),\;(13),\;(23),\;(123),\;(132).
\]
按右侧先做的约定（先做右边的置换），
\[
(12)(23) = (123),\qquad (23)(12) = (132).
\]
二者不等。\(S_3\) 不是交换群。操作顺序不是细节：先交换座位 2 和 3，再交换 1 和 2，与反过来执行，会把每个人送到不同的座位。

置换群与正方形的对称群之间存在深刻的联系。Cayley 定理说：每个群都同构于某个置换群的子群。直觉：让群元素通过左乘去重排群本身。对每个固定的 \(g\in G\)，映射 \(x\mapsto gx\) 是群 \(G\) 到自身的一个双射——一种置换。而且复合关系会完全复现群乘法：\(x\mapsto (gh)x\) 等于先做 \(x\mapsto hx\) 再做 \(y\mapsto gy\)。通过这种方式，任何抽象群都可以被看成某个集合上的可逆重排——群从抽象回到具体操作。

\subsection{矩阵群：可逆线性变换怎样复合}

所有可逆 \(n\times n\) 实矩阵在矩阵乘法下构成\textbf{一般线性群}
\[
GL(n,\R).
\]
单位元是单位矩阵 \(I\)，逆元是通常的矩阵逆。矩阵乘法的结合律来自线性映射复合的结合律；封闭性来自可逆映射的复合仍可逆。行列式为非零实数的全部 \(n\times n\) 矩阵都在这个群里。

正交矩阵组成子群 \(O(n)\)：它们保持内积与长度不变。行列式为 1 的正交矩阵组成\textbf{特殊正交群} \(SO(n)\)，描述保持定向的旋转。平面旋转矩阵：
\[
R(\theta)=\begin{pmatrix}\cos\theta&-\sin\theta\\\sin\theta&\cos\theta\end{pmatrix},
\]
满足
\[
R(\theta)R(\phi)=R(\theta+\phi),\qquad R(\theta)^{-1}=R(-\theta).
\]
角度加法群 \(\R\) 与旋转矩阵群 \(SO(2)\) 之间有一个映射 \(\theta\mapsto R(\theta)\)。它满足 \(R(\theta_1+\theta_2)=R(\theta_1)R(\theta_2)\)——加法的像等于像的乘积。但角度相差 \(2\pi\) 会得到同一个旋转矩阵：\(R(\theta+2\pi)=R(\theta)\)。不同的参数表示同一个操作。这个``多对一''现象后面会格式化为群同态与商群理论——核心问题是：何时不同的表示可以被视为``相同操作的不同记号''，如何用正规子群精确刻画这种等同。

旋转矩阵和刚体运动不止是抽象的例子。在机器人学中，末端执行器的方向由 \(SO(3)\) 中的旋转矩阵描述；两个旋转的复合就是矩阵乘法——永远在群中，永远可逆。在计算机图形学中，场景物体的变换链由若干个平移、旋转、缩放的矩阵乘积实现；只要所有变换矩阵都可逆，整条链就可逆。群的结构让你不必每一步都重新检查可逆性——你知道它还在这扇门里。

矩阵群的价值不只在于举例。物理中的旋转、图网络中的坐标变换、数据增强中的几何变换，都在向同一个抽象提问：哪些可逆线性变换能复合？复合的次序又怎样影响结果？

\subsection{子群：一群中的一小群}

回到正方形的八个对称操作。纯旋转 \(\{e,r,r^2,r^3\}\) 自己也是一个群：任意两个旋转的复合仍是旋转，结合律继承自 \(D_4\)，单位元 \(e\) 在，每个旋转的逆旋转也在。把群公理逐条套上去，全部通过。

这个``大群中的小群''就叫子群。形式定义：\(H\subseteq G\) 是子群，若 \(H\) 自身在 \(G\) 的运算下也构成群。验证子群只需要两条：\(H\) 非空，且对任意 \(a,b\in H\)，有 \(ab^{-1}\in H\)。这两条自动包含单位元在 \(H\) 中和封闭性。

\(D_4\) 中的子群值得清点一遍：
\begin{itemize}
\tightlist
\item 平凡子群 \(\{e\}\)——只有单位元；
\item 旋转子群 \(\{e,r,r^2,r^3\}\)——所有保持定向的操作；
\item 翻折子群 \(\{e,s\}\)、\(\{e,sr\}\)、\(\{e,sr^2\}\)、\(\{e,sr^3\}\)——各含一个翻折和单位元；
\item 由 \(r^2\) 生成的 \(\{e,r^2\}\)——旋转 \(180^\circ\) 两次回到原点；
\item 整个 \(D_4\) 本身。
\end{itemize}

子群的数量和结构不是随机的。Lagrange 定理说：有限群中，子群的阶（元素个数）整除群的阶。\(D_4\) 有 8 个元素，子群的阶只可能是 1、2、4、8，恰好覆盖所有可能。阶为 3 的子群不可能存在——8 不被 3 整除。

子群的概念让``大结构中的小结构''有了精确语言。旋转子群告诉你：即使不知道翻折操作，正方形仍然有一套完整的旋转对称体系。之后轨道-稳定子定理会进一步揭示子群、群作用、对象对称性之间的定量关系。

\subsection{回到正方形：用三条关系压缩一张运算表}

正方形对称群记为 \(D_4\)（有些书叫 \(D_8\)，下标取的是元素个数；这里取二面体群的标准命名）。\(D_4\) 满足三条关系：
\[
r^4=e,\qquad s^2=e,\qquad srs=r^{-1}.
\]

第三条关系是读懂整个群的钥匙。它等价于 \(sr=r^{-1}s\)：先把 \(s\) 移到 \(r\) 左侧，代价是 \(r\) 变成 \(r^{-1}\)。它明确宣告 \(sr\neq rs\)——如果二者相等，代入 \(srs=r^{-1}\) 会推出 \(rss=r^{-1}\)，即 \(r=r^{-1}\)，这与四分之一旋转不等于反向四分之一旋转矛盾。

这三条关系的力量：任意由 \(r\) 和 \(s\) 拼出的长字符串，都可以用它们化简成标准形 \(r^k\) 或 \(sr^k\)（\(k=0,1,2,3\)）。看一个具体的化简过程：
\[
\begin{aligned}
s r^3 s r^2
&= s r^3 (s r) r
   &&\text{(拆开最后的 \(r^2\))}\\
&= s r^3 (r^{-1} s) r
   &&\text{(用 \(sr=r^{-1}s\) 换序)}\\
&= s r^2 s r
   &&\text{(\(r^3 r^{-1}=r^2\))}\\
&= s (r^2 s) r
   &&\\
&= s (s r^{-2}) r
   &&\text{(再次换序: \(r^2s=sr^{-2}\))}\\
&= s^2 r^{-2} r
   &&\\
&= e \cdot r^{-1}
   &&\text{(\(s^2=e\), \(r^{-2}r=r^{-1}\))}\\
&= r^3.
   &&\text{(\(r^{-1}=r^3\))}
\end{aligned}
\]
整个过程只用了三条关系和群公理（结合律、单位元）。不需要看正方形，不需要空间想象。八幅几何图被压缩成三个生成元与关系，之后所有推导都是代数——这就是群的实用价值。

八幅几何图被压缩成两个生成元加三条关系。之后所有组合都可以代数化简，不必再动手旋转纸片。这就是群的核心价值：把大量具体操作压缩成少数规则，然后用这些规则进行严格推理。

\subsection{群 vs 更弱的代数结构：一张对照表}

群的四条公理很强。把其中一条或两条拿掉，掉到的是什么结构？以下是一张速查，每个层级标注了``还可以做什么''和``已经不能做什么''。

\begin{itemize}
\tightlist
\item \textbf{半群 (Semigroup)}：只有封闭性和结合律。例子：正整数在加法下（无单位元，无逆元）。能做的事：反复做同一操作，n 次幂仍有定义。不能做的事：方程 \(a+x=b\) 未必有解，不能撤销。
\item \textbf{幺半群 (Monoid)}：半群 + 单位元。例子：字符串在拼接下（空串为单位元，但没有``反字符串''来撤销）。能做的事：有``什么都不做''的基准。不能做的事：没有撤销，不能解方程。
\item \textbf{群 (Group)}：幺半群 + 每个元素有逆元。能做全部前面结构能做的事，加上：方程 \(ax=b\) 总有唯一解，消去律成立，可逆操作可以串成链。
\item \textbf{Abel 群 (Abelian group)}：群 + 交换律。例子：整数加法。能做群能做的一切事，加上``操作的先后顺序无关紧要''——这个性质在群中已经需要额外验证了。
\end{itemize}

这张对照表的训练价值不在于背诵定义，而在于给你一个筛查流程。遇到一个新的操作集合，不要急着说``这是群''。从半群开始逐级往上筛：封闭性？结合律？有单位元吗？每个元素都能撤销吗？最后才问：操作顺序换一下，结果还一样吗？每通过一级，工具箱就多几个可用的定理。

举个非数学的例子。浏览器的前进/后退按钮：操作集合是浏览历史中的各页面，运算是``前进到下一个''或``后退到上一个''。有单位元（当前页面），有逆元（前进和后退互为逆），但前进再后退到达的页面取决于你当前在哪——这是群在作用，而不是群本身在复合。把它强行塞进群的框架会扭曲原意，把它放在群作用（group action）的框架下才顺理成章。挑选正确的代数结构本身就是一个判断力训练。

\subsection{这一章的准备}

对称操作之所以值得抽象，不只是因为正方形好看。密码学中的置换、物理学中的旋转与 Lorentz 变换、多项式方程根之间的 Galois 置换、数据增强中的平移与翻转、图同构检测中的节点重排——这些不同领域的问题都在追问同一件事：哪些可逆操作能复合？哪些性质在这些操作下保持不变？操作的复合次序是否影响结果？

这一节做的事情是从一个具体的正方形出发，逐步剥离掉几何形状、剥离掉具体的角度和坐标，最后只剩四条公理和一批从公理推导出来的事实。从正方形到抽象群，砍掉了哪些东西？砍掉了操作的具体含义——旋转还是翻折，不重要。只保留了：两个操作可以复合、复合满足结合律、存在什么都不做的操作、每个操作可逆。这些被保留的东西在不同场景下以不同面目出现，但推理规则完全相同。正方形的对称群和经验中学到的``撤销操作是唯一的''或``方程在群里有唯一解''，说的是同一个骨架的不同实例。

下一篇让群真正``作用''到对象上，把``两个东西在对称下看起来相同''精确成轨道与不变量。从正方形纸片的旋转翻折，到抽象的公理和生成元，再回到轨道计数和 Burnside 引理——我们转了一圈，但手中的工具已经完全不一样了。
